\section{Theoretical Proofs}
\label{app:proofs}

This appendix provides the proofs for the theoretical statements in Section~\ref{sec:method}. For a distribution $p$ on a finite set $E$ and a subset $A\subseteq E$, write $p(A)=\sum_{j\in A}p(j)$, and define total variation distance by
\begin{equation}
    \mathrm{TV}_E(p,q)
    =
    \sup_{A\subseteq E}|p(A)-q(A)|
    =
    \frac{1}{2}\sum_{j\in E}|p(j)-q(j)|.
\end{equation}

\subsection{Proof of Lemma 1}

For every subset $A\subseteq C_i$,
\begin{equation}
    |p_i^S(A)-p^T_{i,r}(A)|
    \le
    \mathrm{TV}_{C_i}(p_i^S,p^T_{i,r})
\end{equation}
by the variational definition of total variation distance on $C_i$. Since $p^T_{i,r}$ and $p_i^S$ are both distributions over $C_i$, Pinsker's inequality gives
\begin{equation}
\begin{aligned}
    \mathrm{TV}_{C_i}(p_i^S,p^T_{i,r})
    &\le
    \sqrt{
    \frac{1}{2}
    D_{\mathrm{KL}}\left(p^T_{i,r}\|p_i^S\right)
    }\\
    &\le
    \sqrt{\frac{\delta}{2}}.
\end{aligned}
\end{equation}
Combining the two inequalities proves the claim.

\subsection{Proof of Lemma 2}

On $C_i$, the full diffusion row decomposes as
\begin{equation}
    q_{i,r}(j)=\rho p^T_{i,r}(j),
    \qquad j\in C_i,
\end{equation}
where $\rho=\rho_{i,r}(C_i)$. Outside $C_i$, the zero-extended student distribution has no mass. Therefore
\begin{equation}
\begin{aligned}
&2\,\mathrm{TV}_{\mathcal{D}}
\left(q_{i,r},\widehat p_i^S\right)\\
&\quad=
\sum_{j\in C_i}
\left|\rho p^T_{i,r}(j)-p_i^S(j)\right|
+
\sum_{j\notin C_i}q_{i,r}(j).
\end{aligned}
\end{equation}
For $j\in C_i$,
\begin{equation}
\begin{aligned}
    \left|\rho p^T_{i,r}(j)-p_i^S(j)\right|
    &\le
    \rho\left|p^T_{i,r}(j)-p_i^S(j)\right|\\
    &\quad +(1-\rho)p_i^S(j).
\end{aligned}
\end{equation}
Summing over $j\in C_i$ gives
\begin{equation}
\begin{aligned}
&\sum_{j\in C_i}
\left|\rho p^T_{i,r}(j)-p_i^S(j)\right|\\
&\quad\le
2\rho\,\mathrm{TV}_{C_i}(p^T_{i,r},p_i^S)
+(1-\rho).
\end{aligned}
\end{equation}
Since $\sum_{j\notin C_i}q_{i,r}(j)=1-\rho$, we obtain
\begin{equation}
\begin{aligned}
&2\,\mathrm{TV}_{\mathcal{D}}
\left(q_{i,r},\widehat p_i^S\right)\\
&\quad\le
2\rho\,\mathrm{TV}_{C_i}(p^T_{i,r},p_i^S)
+2(1-\rho).
\end{aligned}
\end{equation}
Dividing by two proves
\begin{equation}
    \mathrm{TV}_{\mathcal{D}}\left(q_{i,r},\widehat p_i^S\right)
    \le
    (1-\rho)
    +
    \rho\,\mathrm{TV}_{C_i}\left(p^T_{i,r},p_i^S\right).
\end{equation}
If $\rho\ge 1-\epsilon$, then $1-\rho\le\epsilon$ and $\rho\le 1$. Applying Lemma 1 to the restricted mismatch yields
\begin{equation}
    \mathrm{TV}_{\mathcal{D}}\left(q_{i,r},\widehat p_i^S\right)
    \le
    \epsilon + \sqrt{\frac{\delta}{2}}.
\end{equation}

\subsection{Proof of Proposition 3}

The mean-square preservation statement follows directly from the definition of $\mathcal{L}_{\mathrm{spec}}(B)$. For any $i,j\in B$, the reverse triangle inequality gives
\begin{equation}
\begin{aligned}
&\left|\left\|z_i^S-z_j^S\right\|_2
-\left\|u_i^T-u_j^T\right\|_2\right|\\
&\quad\le
\left\|(z_i^S-z_j^S)-(u_i^T-u_j^T)\right\|_2.
\end{aligned}
\end{equation}
The triangle inequality then gives
\begin{equation}
\begin{aligned}
&\left\|(z_i^S-z_j^S)-(u_i^T-u_j^T)\right\|_2\\
&\quad\le
\left\|z_i^S-u_i^T\right\|_2
+
\left\|z_j^S-u_j^T\right\|_2.
\end{aligned}
\end{equation}
Averaging over all ordered pairs and applying Jensen's inequality yields
\begin{equation}
\begin{aligned}
&\frac{1}{|B|^2}\sum_{i,j\in B}
\left(\left\|z_i^S-u_i^T\right\|_2
+
\left\|z_j^S-u_j^T\right\|_2\right)\\
&\quad=
\frac{2}{|B|}\sum_{i\in B}
\left\|z_i^S-u_i^T\right\|_2
\le
2\sqrt{\eta}.
\end{aligned}
\end{equation}
This proves the proposition.

\subsection{Proof of Corollary 4}

Let
\begin{equation}
\begin{aligned}
    \mathcal{L}_{\mathrm{HeatGeo}}(B)
    ={}&
    \lambda_{1}\mathcal{L}_{\mathrm{diff}}(B)
    +
    \lambda_{2}\mathcal{L}_{\mathrm{spec}}(B)\\
    &+
    \lambda_{3}\mathcal{L}_{\mathrm{anchor}}(B)
    \le \gamma.
\end{aligned}
\end{equation}
All three losses are nonnegative. Hence
\begin{equation}
\begin{aligned}
    \mathcal{L}_{\mathrm{diff}}(B)
    &\le \frac{\gamma}{\lambda_{1}},\\
    \mathcal{L}_{\mathrm{spec}}(B)
    &\le \frac{\gamma}{\lambda_{2}},\\
    \mathcal{L}_{\mathrm{anchor}}(B)
    &\le \frac{\gamma}{\lambda_{3}}.
\end{aligned}
\end{equation}
Let $\Omega_e=\sum_{r\in\mathcal{R}_e}\omega_r$. Draw $I$ uniformly from $B$ and draw $R$ from $\mathcal{R}_e$ with probability $\Pr(R=r)=\omega_r/\Omega_e$. By the definition of $\mathcal{L}_{\mathrm{diff}}$,
\begin{equation}
\begin{aligned}
&\mathbb{E}_{I,R}\left[
D_{\mathrm{KL}}\left(p^T_{I,R}\|p_I^S\right)
\right]\\
&\quad=
\frac{\mathcal{L}_{\mathrm{diff}}(B)}{\Omega_e}
\le
\frac{\gamma}{\lambda_{1}\Omega_e}.
\end{aligned}
\end{equation}
Applying Lemma 1 and Jensen's inequality gives
\begin{equation}
\begin{aligned}
&\mathbb{E}_{I,R}\left[\sup_{A\subseteq C_I}
\left|p_I^S(A)-p^T_{I,R}(A)\right|\right]\\
&\quad\le
\sqrt{
\frac{\gamma}{2\lambda_{1}\Omega_e}
}.
\end{aligned}
\end{equation}
Combining this bound with Lemma 2 and the assumption $\rho_{i,r}(C_i)\ge 1-\epsilon$ gives
\begin{equation}
\begin{aligned}
&\mathbb{E}_{I,R}\left[
\mathrm{TV}_{\mathcal{D}}
\left(q_{I,R},\widehat p_I^S\right)
\right]\\
&\quad\le
\epsilon+
\sqrt{
\frac{\gamma}{2\lambda_{1}\Omega_e}
}.
\end{aligned}
\end{equation}
Proposition 3 with $\eta=\gamma/\lambda_{2}$ gives
\begin{equation}
\begin{aligned}
&\frac{1}{|B|^2}\sum_{i,j\in B}
\left|\left\|z_i^S-z_j^S\right\|_2
-\left\|u_i^T-u_j^T\right\|_2\right|\\
&\quad\le
2\sqrt{\frac{\gamma}{\lambda_{2}}},
\end{aligned}
\end{equation}
and the anchor bound follows directly:
\begin{equation}
    \frac{1}{|B|}\sum_{i\in B}
    \left(1-\cos(W_as_i,t_i)\right)
    \le
    \frac{\gamma}{\lambda_{3}}.
\end{equation}
